% ============================================================
%  harryopo-book 书籍模板示例（ctexbook 骨架 + 书籍化结构）
%  用法: office.py render input.md --format book
%        office.py render input.md --template harryopo-book
%  本文件是注册表登记用的展示样张：手册写 .tex 里定义的命令与结构，
%  实际渲染走 convert.py --type book（MD → TEX → xelatex ×3）。
% ============================================================
\documentclass{harryopo-book}

\title{书名：示例手册}
\subtitle{副标题：harryopo 书籍模板能力演示}
\author{作者名}
\bookseries{丛书名}
\institute{机构名}
\date{2026 年 9 月 30 日}

\begin{document}

% ===================== frontmatter =====================
\frontmatter
\maketitle
\tableofcontents

% ===================== mainmatter =====================
\mainmatter

\bookabstract{
本书演示 harryopo-book 文档类的能力：五级方正字体阶梯、章级编号的定理
环境、六种教学盒子、右侧旁注与双面页眉。
}

\chapter{章节结构}
\section{节与五级字体阶梯}
\subsection{小节：黑体}
\subsubsection{小小节：楷体}
\subhead{括号式小标题}
正文用方正书宋；加粗用 {\fzht 黑体} 而非 \textbf。

% 右侧旁注（marginparwidth=4cm）
\bianzhu[旁注]{这是右侧旁注，仿宋小字，用于补充说明或给出参考链接。}

% 定理环境：章级编号，且定义/命题/引理/推论共享计数器
\begin{theorem}[勾股定理]
直角三角形两直角边平方和等于斜边平方。
\end{theorem}
\begin{definition}
三条边分别称为直角边与斜边。
\end{definition}
\begin{remark}
本环境与定义共享编号序列。
\end{remark}

\begin{equation}
  \int_0^1 x^2 \,\mathrm{d}x = \frac{1}{3}
\end{equation}

\chapter{六种教学盒子}
\section{核心概念}
\begin{keyconcept}
向量空间是线性代数的基本对象。
\end{keyconcept}

\section{例}
\begin{examplebox}
求矩阵 $A$ 的特征值。
\end{examplebox}

\section{小贴士}
\begin{tipbox}
计算前先看清题目要的是行空间还是列空间。
\end{tipbox}

\section{注意}
\begin{warningbox}
相似不等于合同；行列式与迹在相似变换下不变。
\end{warningbox}

\section{练习}
\begin{practicebox}
证明：可逆矩阵的逆矩阵可逆。
\end{practicebox}

\section{解答}
\begin{answerbox}
设 $A^{-1}$ 存在，则其行列式与 $A$ 互为倒数，故均非零。
\end{answerbox}

% ===================== backmatter =====================
\backmatter
\chapter*{参考文献}
\addcontentsline{toc}{chapter}{参考文献}
\begin{thebibliography}{99}
\bibitem{ref1} Knuth D E. The TeXbook.
\bibitem{ref2} Lamport L. LaTeX: A Document Preparation System.
\end{thebibliography}

\end{document}